\documentclass[10pt,a4paper]{article}
\usepackage{graphicx} % Required for inserting images
\usepackage{setspace}
\usepackage[margin=2.5cm]{geometry}
\usepackage{amsmath, amssymb, amsfonts}
\usepackage{mdframed}
\usepackage{float}
\newmdenv[
  linewidth=0.7pt,
  roundcorner=3pt,
  skipabove=6pt,
  skipbelow=6pt,
  innerleftmargin=6pt,
  innerrightmargin=6pt,
  innertopmargin=4pt,
  innerbottommargin=4pt
]{algobox}
\usepackage{enumitem}

\title{BIS638 Final Project}
\author{Manchen Wang}
\date{November 2025}

\begin{document}
\doublespacing
\maketitle

\section{Introduction}

\section{Data Manipulation and Visualization}
\subsection{Entity Relationship Diagram}

As shown in Figure 1 below, the ERD provides a high-level overview of how the major MIMIC-IV tables are connected in our cohort construction process. Rather than reproducing each technical SQL join, the diagram highlights the conceptual data flow that supports our pneumonia-to-ICU study. We begin by identifying patients with pneumonia using ICD-10 codes from the Diagnoses ICD table. These diagnosis events link to the Admissions table through hadm id, allowing us to retrieve demographic information and in-hospital outcomes such as mortality. Admissions are further connected to the ICU Stay Detail table through stay id, which enables us to isolate ICU encounters for the pneumonia cohort. Once an ICU stay is defined, we attach first-day clinical measurements—including labs, vital signs, blood gas values, urine output, height, and weight—by joining their respective tables using stay id.
This ERD serves as a roadmap that helps the reader understand how heterogeneous clinical data sources combine into one analytic dataset. By focusing on structural relationships rather than code-level details, the diagram offers a clear and intuitive summary of how our analytic cohort is assembled for demographic analysis and predictive modeling.


\subsection{ Ontology Graph}

We also build a pneumonia ontology graph that provides a more holistic, scientific view of the disease beyond the structured variables in MIMIC-IV. While the ERD focuses on how data tables connect, the ontology graph highlights how clinical concepts related to pneumonia connect to one another. It synthesizes knowledge from medical literature and standard biomedical ontologies to show causes, risk factors, symptoms, treatments, and potential outcomes. By organizing these relationships visually—for example, linking pneumonia to pathogens, clinical signs, mechanical ventilation, or downstream risks such as ICU admission and mortality—the graph helps us build intuitive domain understanding before modeling.

\begin{figure}[H]
\centering
\begin{minipage}{0.48\linewidth}
    \centering
    \includegraphics[width=\linewidth]{Screenshot 2025-12-05 at 09.33.56.png}
    \caption{Entity–Relationship Diagram}
    \label{fig:erd}
\end{minipage}
\hfill
\begin{minipage}{0.48\linewidth}
    \centering
    \includegraphics[width=\linewidth]{Screenshot 2025-12-05 at 09.39.48.png}
    \caption{Pneumonia Ontology Graph}
    \label{fig:ontology}
\end{minipage}
\end{figure}






\section{Dashboard}

\subsection{Method}
We developed an interactive web-based dashboard to support exploratory analysis of pneumonia patient outcomes using the MIMIC-IV database. The dashboard is a single-page web application built in React~18 with TypeScript, bundled with Vite and styled with Tailwind CSS. A CSV file containing the patient data is loaded and processed locally using \texttt{PapaParse} and stored in global state via the React Context API. Users filter records by age, gender, race, insurance, admission type, outcome, and length of stay using sliders, multi-select dropdowns, and search boxes. The filtering logic applies all active criteria and passes the resulting subset to the chart components. Visualizations are implemented with \texttt{Recharts} (bar charts, histograms, pie charts, and area plots). This dashboard enables flexible exploration of patient data without database queries. The system architecture in Figure~\ref{fig:architecture} illustrates this workflow.

\begin{figure}[H]
    \centering
    \includegraphics[width=0.6\textwidth]{architecture_diagram.jpg}
    \caption{Architecture}
    \label{fig:architecture}
\end{figure}

\subsection{Results}

The dashboard contains eight different specialized pages. Here, we present three representative examples (Figures~\ref{fig:overview}--\ref{fig:vitals}) to demonstrate key features.

\textbf{Overview Analytics.} Figure~\ref{fig:overview} displays the summary interface with key performance indicators presented as metric cards, including total patient counts, mortality rates, and average length of stay measurements. Multiple chart types (donut charts, stacked bar charts, and area plots) illustrate outcome distributions, age demographics, and ICU stay patterns stratified by survival status.

\textbf{Demographic Visualization.} The demographics module (Figure~\ref{fig:demographics}) employs stacked bar charts to visualize patient stratification across multiple dimensions: gender, age groups, race/ethnicity, insurance type, and marital status. Each visualization compares survival outcomes within demographic categories, facilitating health disparity analysis and cohort comparisons.

\textbf{Vital Signs Analysis.} Figure~\ref{fig:vitals} demonstrates the vital signs module with tabbed navigation across seven physiological parameters (heart rate, blood pressure, respiratory rate, temperature, SpO2). For each vital sign, the interface displays statistical summaries (minimum, mean, maximum values) comparing survived versus deceased cohorts, accompanied by distribution histograms. Clinical normal ranges are annotated for reference.

\begin{figure}[H]
    \centering
    \includegraphics[width=0.8\textwidth]{overview.png}
    \caption{Dashboard Overview displaying KPIs, outcome distribution donut chart, age distribution histogram, and ICU length of stay area chart.}
    \label{fig:overview}
\end{figure}

\begin{figure}[H]
    \centering
    \includegraphics[width=0.8\textwidth]{demographics.png}
    \caption{Demographics page showing gender distribution, age-stratified survival analysis, and race/ethnicity breakdown with mortality rates.}
    \label{fig:demographics}
\end{figure}

\begin{figure}[H]
    \centering
    \includegraphics[width=0.8\textwidth]{vital_signs.png}
    \caption{Vital signs analysis page displaying heart rate statistics for survived and deceased patients, with distribution histogram and normal range reference.}
    \label{fig:vitals}
\end{figure}

\section{Machine Learning Analysis of ICU Pneumonia Mortality}



\section{Bayesian Hierarchical Logistic Model}
\subsection{Model Specification and Algorithm}
Our ultimate goal is to improve the model interpretability by probabilistic variable selection. To achieve this, we need a suitable statistical model. However, simply applying the ordinary logistic regression to the binary outcome is not enough, given the structure of our dataset. Many patients have multiple ICU stay records; these records are correlated rather than independent and identically distributed due to the patient-specific effect. This indicates the need to apply the Bayesian hierarchical logistic model.\\
We let $i \in \{1,\dots,n\}$ be the indices for the ICU stay IDs, $j\in \{1,\dots,n_{\text{id}}\}$ be the indices for the subject IDs, and $id_i$ be the subject ID for the i-th ICU stay record. We added both patient-level and record-level randomness, and denoted $b_j$ as the patient level random intercept. We also apply a binary indicator $z_k$, $k \in \{1, \dots, p\}$ as indices for features, for feature selection. Then the new model would be
\[
y_i \mid p_i, b_{\mathrm{id}_i} \sim \mathrm{Bernoulli}(p_i).
\]
where
\[
\text{logit}(p_i)
= x_i^\top (z\odot \beta) + b_{\mathrm{id}_i},\quad\text{and}\quad b_j \sim \mathrm{N}(0,\sigma_b^2),\, j=1,\dots,n_{\text{id}}.
\]
This yields a hierarchical logistic regression model, as the parameter $b_{\mathrm{id}_i}$ modelling $y_i$ depends on another parameter $\sigma_b$. More specifically, the model could be divided into two "submodels", with $\beta$ and $b$ having likelihood constructed by y, and $\sigma_b$ having likelihood constructed by $b$.\\
We adopted the weakly informative prior, $\beta_k \sim N(0, 100)$, $z_k \sim \text{Bern}(0.5)$ and $\sigma_b^2 \sim \text{IG}(0.001, 0.001)$. Under such model, the analytical form of the posterior is intractable. Therefore, we implemented the Gibbs sampling to get the posterior samples. For the logistic setting, we introduce a efficient Gibbs sampling method, called the P\'olya-Gamma Data Augmentation method. To accommodate the sampling process, we construct an ID matrix, $Z \in \mathbb{R}^{n\times n_{\mathrm{id}}},\, Z_{i,k} = \mathbf{1}\{id_i = k\}$, and stack it with $X$ to construct a new covariate matrix, $\tilde X = \bigl[X\; Z\bigr] \in \mathbb{R}^{n\times (p+n_{\mathrm{id}})}$. Then, if we define $\theta = (\beta^\top,b^\top)^\top$, our logit function could be written as $\text{logit}(p_i) = \tilde{x}_i^\top \theta$,  and the prior for $\theta$ is $\theta \sim N\bigl(0, \mathrm{diag}\!\bigl(100 I_p,\,\sigma_b^2 I_{n_{\mathrm{id}}}\bigr)\bigr)$. The resulting Gibbs sampling procedure is as follows:
\begin{algobox}
\begingroup
\small
\begin{spacing}{0.9}

\noindent\textbf{1.\ Update $(\omega,\beta,b)$ by P\'olya-Gamma:}
\begin{enumerate}[itemsep=2pt, topsep=2pt, leftmargin=*]
  \item Set the current values
        $\beta \leftarrow \beta^{(s)}$,
        $z \leftarrow z^{(s)}$,
        $b \leftarrow b^{(s)}$ and
        $\sigma_b^2 \leftarrow \bigl(\sigma_b^2\bigr)^{(s)}$.
  \item Compute the \textbf{linear predictors} $\eta_i
      = x_i^\top (z \odot \beta) + b_{id_i},
      \, i=1,\dots,n,
    $.
  \item Draw the \textbf{latent P\'olya-Gamma variables} $\omega_i \mid \beta,z,b,y \sim PG(1, \eta_i), \, i=1,\dots,n.$
  \item Let $\Omega = \mathrm{diag}(\omega_1,\dots,\omega_n)$, $\kappa = (y_1 - 1/2,\dots,y_n - 1/2)^\top$ and $\Sigma_0^{-1}= \mathrm{diag}\!\bigl(
               10^{-2} I_p,\,
               (\sigma_b^2)^{-1} I_{n_{\mathrm{id}}}
            \bigr)$.
        Compute the \textbf{full conditional covariance and mean} $V_\omega
          = \bigl(\tilde X^\top \Omega \tilde X + \Sigma_0^{-1}\bigr)^{-1},
          \,
          m_\omega
          = V_\omega \tilde X^\top \kappa$. Draw
    \[
      \theta^{(s+1)} = 
      \begin{pmatrix}
        \beta^{(s+1)} \\
        b^{(s+1)}
      \end{pmatrix}
      \sim N\bigl(m_\omega, V_\omega\bigr).
    \]
\end{enumerate}

\medskip
\noindent\textbf{2.\ Update $\sigma_b^2$:}
\begin{enumerate}[itemsep=2pt, topsep=2pt, leftmargin=*]
  \item Given the inverse-gamma prior
        $\sigma_b^2 \sim \mathrm{IG}(a_0,b_0)$ and
        the current $b^{(s+1)}$, sample
        \[
          \bigl(\sigma_b^2\bigr)^{(s+1)}
          \sim \mathrm{IG}\!\left(
                 a_0 + \frac{n_{\mathrm{id}}}{2},\,
                 b_0 + \frac12 \sum_{k=1}^{n_{\mathrm{id}}}
                               \bigl(b_k^{(s+1)}\bigr)^2
               \right).
        \]
\end{enumerate}

\medskip
\noindent\textbf{3.\ Update $z$:}
\begin{enumerate}[itemsep=2pt, topsep=2pt, leftmargin=*]
  \item Set the current value $z \leftarrow z^{(s)}$.
  \item For $j=1,\dots,p$ in random order:
  \begin{enumerate}[itemsep=1pt, topsep=1pt, leftmargin=*]
    \item Compute the Bernoulli probability with the current value of $z$
      \begin{align*}
        p_j
        &= \Pr\bigl(z_j=1 \mid \beta^{(s+1)},b^{(s+1)}, z_{-j}, y, X, id\bigr)\\
        &= \frac{p\bigl(y \mid \beta^{(s+1)}, b^{(s+1)},
                           z_j=1, z_{-j}, X, id\bigr)}
               {p\bigl(y \mid \beta^{(s+1)}, b^{(s+1)},
                           z_j=1, z_{-j}, X, id\bigr)
                + p\bigl(y \mid \beta^{(s+1)}, b^{(s+1)},
                           z_j=0, z_{-j}, X, id\bigr)},
      \end{align*}
    \item Draw $z_j \sim \mathrm{Bernoulli}(p_j)$ and replace the $j$th component of $z$ with this new value.
  \end{enumerate}
\end{enumerate}

\end{spacing}
\endgroup
\end{algobox}
We applied four multiple chains for sampling. For each chain, we discarded the 2,000 burn-in period and performed 8,000 further iterations (a total of 32,000 samples). To measure the quality of the samples, we computed the Effective Sample Size, which estimates the independent draws from all the samples. The median among all feature coefficients, $\beta_j$, is 23,641, showing the samples are well mixing.
\subsection{Feature Selection via Posterior Inclusion Probability}
To quantify the feature importance by probability, we introduced the posterior inclusion probability (PIP). The PIP for the k-th feature is defined as $PIP_k = \mathbb{E}(z_k \mid y)$. By Monte Carlo method, we can apply the samples to estimate the expectation $PIP_k \approx \frac{1}{M}\sum_{s=1}^Mz_s$, where $M$ is the total number of the generated samples.\\
Fig \ref{fig:modelcomp} shows the selected features of the Bayesian model based on the threshold $PIP>0.5$, along with the top features of the machine learning models. As shown in the figure, there is a significant overlap between the Bayesian and the Machine Learning models. Some of the universally recognized key factors, such as \texttt{temperature}, \texttt{urine output}, and \texttt{heart rate}, are selected by all three models, confirming that our Bayesian model correctly identifies critical signals in the data.\\
Meanwhile, some non-overlapping features still exist. There are two main reasons for the distinction. First, our Bayesian model is based on logistic regression, so it tends to select only stable linear features, whereas our Machine Learning models are non-linear and is likely to highlight both linear and non-linear features. Secondly, the Machine Learning models tend to treat each ICU stay record independently, but our Bayesian model is hierarchical. While explicitly accounting for the patient-level random intercept, the Bayesian model is more likely to accept features with no patient-specific variation, thereby being more constrained.
\begin{figure}[htbp]
  \centering
  \includegraphics[width=1\textwidth]{model comparison.png} 
  \caption{The comparison between the selected features from the Bayesian model and the top features from the Random Forest and the XGBoost model. Each row of the colored dots shows the selected features from the corresponding model.}
  \label{fig:modelcomp}
\end{figure}
\subsection{Posterior Predictive Check}
Based on the selected features, we also checked the model accuracy. For a given set of posterior samples, $\beta, z, \sigma_b^2$, since each subject from the test set is unseen, we first sample its random effect by $N(0, \sigma_b^2)$. Then, for the i-th record, we predict the Bernoulli probability for the outcome $y_i$, $p_i = \text{sigmoid}(x_i^\top (z\odot \beta) + b_{\mathrm{id}_i})$. The final predictive Bernoulli probability is the average of the predictive probabilities across all posterior samples.\\
Comparing with the actual outcome in the test set, the model achieved an AUC of around 0.71. Although the accuracy is not very high, it is a reasonable and robust value in this Bayesian setting. One of the reasons is similar to the previous non-overlapping feature interpretation. The Bayesian one imposes more constraints, as it tends to accept only stable linear features with no patient-specific effects. Moreover, as the train-test split is at the patient level, we have no previous knowledge of the records in the test set. Therefore, we assign a completely random intercept to each test subject only based on the sampled individual-level variability, which naturally limits some accuracy.
\section{Conclusion and Discussion}













The Bayesian hierarchical logistic model provides a robust, interpretable model for feature selection while filtering out patient-specific correlations. Looking forward, we could further improve the model accuracy while preserving interpretability by adding nonlinear components, such as Bayesian neural networks and Gaussian Processes, into the Bayesian hierarchical model. Meanwhile, our current evaluation was limited to cases where we had a completely new patient as the test subject. However, in practice, when we are dealing with Electronic Health Records (EHR), longitudinal data will store previous information on the patient-specific effect, and the patient-specific random intercept is dynamically updated to improve prediction accuracy.\\
Based on our current analysis, we suggest putting the overlapping features, such as the age, heart rate, temperature, urine output, albumin and bilirubin at high test priority. Therefore, clinicians can effectively assess the severity of ICU pneumonia. This provides early warning and treatment before deterioration happens. The evaluation also provides suggestions for allocating hospitalization resources. After appropriate interventions are taken, ICU mortality from pneumonia could be better controlled.
\end{document}
